% Modern editorial, markup and figure/layout contributions: CC-BY-SA-4.0.
% Historical text and illustrations: Leonhard Euler (1744), public domain.
% Component scope and upstream attribution: see RIGHTS.md.
\sourcepage{2}
\subsection*{Corollarium II.}
\originalsection{3}Cum autem eadem curva infinitis modis sui similis effici queat, Problema, nisi quædam restrictio adhibeatur, maxime esset indeterminatum, atque adeo nullum. Quæcunque enim curva præbeatur maximi minimive proprietate prædita, semper alia, illi quidem vel similis vel dissimilis, exhiberi posset, quæ illam proprietatem, vel majorem, vel minorem, in se contineret.

\subsection*{Corollarium III.}
\originalsection{4}Quoniam igitur adæquata curvarum cognitio postulat, ut eæ ad axem aliquem positione datum, ejusque portiones quascunque quæ abscissæ vocantur, referantur: prima eaque præcipua restrictio ex quantitate abscissæ petenda erit.

\subsection*{Corollarium IV.}
\originalsection{5}Problemata ergo ad methodum hanc pertinentia ita proponi debent, ut quærantur lineæ curvæ ad axem positione datum relatæ, quæ inter omnes alias curvas eidem abscissæ respondentes maximi minimive proprietate sint præditæ.

\subsection*{Scholion.}
\originalsection{6}Hæc itaque Methodus maximorum \& minimorum maxime discrepat ab illa, quam alibi exposuimus. Ibi enim, pro data ac determinata linea curva, locum determinavimus, ubi proposita quædam quantitas variabilis ad curvam pertinens fiat maxima vel minima. Hic autem ipsa linea curva quæritur, in qua quantitas quædam proposita fiat maxima vel minima. Methodus hæc jam superiori Seculo, mox post inventam Analysin infinitorum, excoli cœpit a Celeb. Fratribus BERNOULLIIS, atque ex eo tempore maxima cepit incrementa. Primum quidem Problema, quod ex hoc genere est tractatum, ad Mechanicam respiciebat, eoque quærebatur linea curva super qua grave descendens citissime delabatur; cui \textit{Curvæ brachystochronæ seu Lineæ celerrimi descensus} nomen erat impositum. In hoc Problemate
\sourcepage{3}
jam manifestum est, id, sine adjuncta conditione, nequidem nomen quæstionis retinere posse: perspicuum enim est, quo brevior magisque ad situm verticalem accedens linea capiatur, eo fore tempus descensus super ea brevius. Quamobrem non absolute quæri potest linea, super qua grave descendens celerrime seu brevissimo tempore delabatur; sed abscissæ quantitas, cui curva invenienda respondeat, simul debuit definiri; ita ut, inter omnes curvas eidem abscissæ in axe positione dato sumtæ respondentes, quæreretur ea super qua corpus grave citissime delaberetur. Neque vero in hoc Problemate ista conditio sufficiebat ad id determinatum efficiendum: sed insuper istam conditionem adjicere oportuit, ut curva invenienda per data duo puncta transeat; atque istud Problema his conditionibus adstringi debuit, ut fieret determinatum, inter omnes, scilicet, lineas curvas per data duo puncta transeuntes eam determinare super qua corpus descendens arcum datæ abscissæ respondentem brevissimo tempore absolvat. Interim tamen hic notandum est, conditionem transitus per duo puncta non esse absolute necessariam, sed in hoc Problemate per ipsam solutionem esse illatam. In solutione enim hujus Problematis immediate pervenitur ad æquationem differentialem secundi gradus, quæ bis integrata duas recipit constantes arbitrarias, ad quas determinandas duobus opus est punctis per quæ curva traducatur, vel aliis similibus proprietatibus: atque hæc eadem conditio, quasi sua sponte, ad omnia istiusmodi Problemata accedit, quarum solutio immediate ad æquationem differentialem secundi gradus deducit. In Problematibus autem quæ resolvuntur per æquationem differentialem quarti vel altioris ordinis, nequidem duo puncta ad curvam determinandam sufficiunt, sed tot opus est punctis, quot gradus differentialia obtinent. Contra vero, si solutio statim ad æquationem algebraicam perducat, tum sine hujusmodi conditione Problema perfecte erit determinatum; dummodo abscissæ longitudo definiatur. Verum hæc omnia clarius perspicientur, quando infra ad solutiones Problematum perveniemus: ibique has notationes fusius explicabimus. Hic enim in principio ista tantum commemorare visum est, ut
\sourcepage{4}
perversas ideas circa determinationem hujusmodi Problematum tollamus.

\subsection*{Definitio II.}
\originalsection{7}\textit{Methodus maximorum ac minimorum absoluta}, docet inter omnes omnino curvas, ad eandem abscissam relatas, determinare eam, in qua proposita quædam quantitas variabilis maximum minimumve obtineat valorem.

\subsection*{Corollarium.}
\originalsection{8}In Problematibus igitur ad hanc methodum pertinentibus, datur axis positione; atque, inter omnes curvas quæ ad hunc axem ejusque determinatam portionem referri possunt, determinatur ea in qua quantitas quædam variabilis sit maxima vel minima.

\subsection*{Scholion.}
\originalsection{9}Aliam conditionem ad maximi minimive determinationem, præter abscissæ quantitatem, hic in genere non adjicimus. Dantur enim Problemata, quæ hoc modo perfecte determinantur; quemadmodum infra distinctius patebit. Etsi enim etiam ejusmodi Problemata occurrunt, ad quæ determinanda insuper duo plurave puncta præscribi possunt, per quæ quæsita curva transeat; tamen hoc demum ex ipsa cujuscunque Problematis solutione perspicietur. Namque si ad ejusmodi æquationem pro curva quæsita perveniatur, in qua per integrationem novæ quantitates constantes sint ingressæ, quæ in ipsa quæstione non inerant; tum solutio censenda erit ambigua, atque vaga; eo quod innumerabiles lineas curvas, quæ ex determinatione illarum quantitatum constantium \& arbitrariarum oriri possunt, in se complectitur. His igitur in casibus erit concludendum, Problema ex sua natura non penitus esse determinatum: sed ad ejus plenam determinationem, præter abscissæ quantitatem, tot novas conditiones adjungi oportere, quibus illæ arbitrariæ constantes ad determinatos valores revocentur. Pro hujusmodi autem conditionibus commodissime assumuntur puncta, per quæ curvæ
